\subsection{Envelope Constriction}

\subsubsection{CSA Two-Stage Calibration}
\begin{itemize}
    \item Training vectors are split into two disjoint halves $S_1$ and $S_2$ (50/50 random split).
    \item $S_1$ (shape discovery): used to learn the geometry of the envelope.
    \item $S_2$ (size scaling): used to calibrate the scaling factor $\hat{t}$ that determines the final size of the envelope.
    \item The two stages must use disjoint data to preserve the exchangeability argument underlying the coverage guarantee.
\end{itemize}

\subsubsection{$\tau$ Computation}
\begin{itemize}
    \item Define $\tau(\mathbf{s}) = \inf\{t > 0 : \mathbf{s} \in t \cdot E\}$, the minimum scaling factor needed to include point $\mathbf{s}$ in the scaled envelope $E$.
    \item The global scaling factor is then determined by the $(1-\alpha)$ quantile of the $\tau$ values over $S_2$:
    \[
    \hat{t} = \tau_{(k)} \quad \text{where} \quad k = \lceil(|S_2|+1)(1-\alpha)\rceil
    \]
    where $\tau_{(k)}$ denotes the $k$-th smallest value of $\tau$ in $S_2$.
    \item Key contribution: for each of the three envelope geometries, $\tau(\mathbf{s})$ has a closed-form expression computable in a single pass — no iterative search required.
    \item Computational benefit: the entire calibration step reduces to one forward pass over $S_2$ followed by a sort.
    \item Contrast with the original CSA paper which uses iterative binary search.
\end{itemize}

\subsubsection{Radial Method}

\subsubsection{Strip Method}

\subsubsection{1D}

\subsubsection{Class-Specific $\hat{t}$ Calculation}


3.4.1 — CSA Two-Stage Calibration

% Training NC vectors split 50/50 into S1S_1
% S1​ and S2S_2
% S2​
% S1S_1
% S1​: shape discovery — learn the geometry of the envelope
% S2S_2
% S2​: size scaling — calibrate t^\hat{t}
% t^
% Disjoint split required to preserve exchangeability

% *3.4.2 — Analytical τ\tau
% τ Computation* **(original contribution)**
% Define τ(s)=inf⁡{t>0:s∈t⋅E}\tau(\mathbf{s}) = \inf\{t > 0 : \mathbf{s} \in t \cdot E\}
% τ(s)=inf{t>0:s∈t⋅E}
% Global scaling factor: t^=τ⌈(∣S2∣+1)(1−α)⌉\hat{t} = \tau_{\lceil(|S_2|+1)(1-\alpha)\rceil}
% t^=τ⌈(∣S2​∣+1)(1−α)⌉​, the (1−α)(1-\alpha)
% (1−α) quantile of τ\tau
% τ-values over S2S_2
% S2​
% Key contribution: for each of the three envelope geometries, τ(s)\tau(\mathbf{s})
% τ(s) has a closed-form expression — no iterative search required
% Computational benefit: entire calibration reduces to one forward pass over S2S_2
% S2​ followed by a sort
% Explicit contrast with the original CSA paper which uses iterative binary search

% 3.4.3 — Radial Method

% Sample M=200M = 200
% M=200 unit directions {um}\{\mathbf{u}_m\}
% {um​} on positive orthant of unit sphere in RK\mathbb{R}^K
% RK
% Per-sector quantile q~m\tilde{q}_m
% q~​m​: (1−α)(1-\alpha)
% (1−α) quantile of magnitudes of S1S_1
% S1​ points within δ=30°\delta = 30°
% δ=30° of um\mathbf{u}_m
% um​; global fallback if fewer than 5 points
% τ\tau
% τ formula: τ(s)=min⁡mmax⁡kskq~m⋅umk\tau(\mathbf{s}) = \min_m \max_k \frac{s_k}{\tilde{q}_m \cdot u_{mk}}
% τ(s)=minm​maxk​q~​m​⋅umk​sk​​ over sectors covering s\mathbf{s}
% s
% Test-time membership: ∃m\exists m
% ∃m such that sk≤t^⋅q~m⋅umks_k \leq \hat{t} \cdot \tilde{q}_m \cdot u_{mk}
% sk​≤t^⋅q~​m​⋅umk​ for all kk
% % k
% 3.4.4 — Strip Method

% Data-adaptive bin edges: M=20M = 20
% M=20 equal-width bins per dimension using min/max of S1S_1
% S1​
% Per-(j,i,m)(j, i, m)
% (j,i,m) quantile: (1−α)(1-\alpha)
% (1−α) quantile of dimension ii
% i among S1S_1
% S1​ points in bin mm
% m of dimension jj
% j; global fallback if fewer than 3 points
% Monotonicity propagation: sweep right to left, enforce limits[j,i,m]≥limits[j,i,m+1]\text{limits}[j,i,m] \geq \text{limits}[j,i,m+1]
% limits[j,i,m]≥limits[j,i,m+1]
% τ\tau
% τ formula: τ(s)=max⁡j≠isilimits[j,i,binj(s)]\tau(\mathbf{s}) = \max_{j \neq i} \frac{s_i}{\text{limits}[j, i, \text{bin}_j(\mathbf{s})]}
% τ(s)=maxj=i​limits[j,i,binj​(s)]si​​
% Test-time membership: ∀j,∀i≠j:si≤t^⋅limits[j,i,binj(s)]\forall j, \forall i \neq j: s_i \leq \hat{t} \cdot \text{limits}[j, i, \text{bin}_j(\mathbf{s})]
% ∀j,∀i=j:si​≤t^⋅limits[j,i,binj​(s)]
% Why it is non-convex: the allowed range of dimension ii
% i varies depending on where along dimension jj
% j the point sits
% 3.4.5 — Collapsed 1D Method

% Reduce s\mathbf{s}
% s to scalar: sˉ=1K∑ksk\bar{s} = \frac{1}{K}\sum_k s_k
% sˉ=K1​∑k​sk​
% q~\tilde{q}
% q~​: (1−α)(1-\alpha)
% (1−α) quantile of S1S_1
% S1​ scalars
% τ\tau
% τ formula: τ(s)=sˉ/q~\tau(\mathbf{s}) = \bar{s} / \tilde{q}
% τ(s)=sˉ/q~​
% Test-time membership: sˉ≤t^⋅q~\bar{s} \leq \hat{t} \cdot \tilde{q}
% sˉ≤t^⋅q~​
% Trade-off: simplest method, unaffected by curse of dimensionality, but discards all geometric structure

% *3.4.6 — Class-Conditional t^\hat{t}
% t^ Calibration*

% Problem: single t^\hat{t}
% t^ gives marginal coverage, can under-cover minority classes
% Fix: compute t^y\hat{t}_y
% t^y​ separately per class from S2S_2
% S2​ points with label yy
% y, set t^=max⁡yt^y\hat{t} = \max_y \hat{t}_y
% t^=maxy​t^y​
% Formal statement: this ensures P(Y∈C(X)∣Y=y)≥1−αP(Y \in C(X) \mid Y = y) \geq 1 - \alpha
% P(Y∈C(X)∣Y=y)≥1−α for each class individually
% Computational cost: negligible — per-class sort of already-computed τ\tau
% τ values
% Formula: t^y=τ⌈(∣S2y∣+1)(1−α)⌉\hat{t}_y = \tau_{\lceil(|S_2^y|+1)(1-\alpha)\rceil}
% t^y​=τ⌈(∣S2y​∣+1)(1−α)⌉​ where S2y={s∈S2:Y=y}S_2^y = \{\mathbf{s} \in S_2 : Y = y\}
% S2y​={s∈S2​:Y=y}